% Source: https://artofproblemsolving.com/community/c3h2616144_jmpsc_sprint_problem_19
% Stable identity: aops-topic-2616144-q01
The hour digit changes from $M$ to $A\ne M$, so the hour advanced: $A=M+1$, and the minutes went past $60$. So
\[\overline{AT}+14-60=\overline{TM}\qquad\Longrightarrow\qquad10A+T-46=10T+M.\]
Substituting $A=M+1$ gives $9M-36=9T$, so $T=M-4$. Also $\overline{AT}\ge46$, so the minutes' tens digit is $A=4$ or $5$. $A=4$ gives $M=3$ and $T=-1$, which is impossible. $A=5$ gives $M=4$ and $T=0$: the times are $4{:}50$ and $5{:}04$, which fit. So $\textbf{A}\times\textbf{M}=5\cdot4=\boxed{20}$.
